\section{Experimental Results}

\subsection{Round-1 Results}
Table~\ref{tab:main-results} reports the currently available means over the 15
severity-5 corruption streams. On CIFAR-10-C with noise outliers, KNN--STAMP
obtains the highest known-class accuracy, improving on the frozen source model
by 21.30 points and on test-time normalization by 5.69 points. This gain does
not transfer to outlier detection: its AUC is 9.41 points below normalization,
and its H-score is 6.97 points lower. On CIFAR-100-C, KNN--STAMP improves all
three metrics over the source model, but is approximately tied with
normalization in accuracy ($-0.17$ points) while trailing it by 10.98 AUC points
and 4.34 H-score points.

These results neither show a uniform failure nor support the original
hypothesis. The adaptation machinery remains useful for recognition, especially
on CIFAR-10-C, but feature distance produces a poorer joint trade-off than the
simple normalization baseline. Aggregate metrics cannot establish the cause;
in particular, they do not reveal whether KNN rejects corrupted ID samples or
admits noise outliers. The admission diagnostics added for round 2 are intended
to distinguish these failure modes.

\begin{table}[t]
    \caption{Mean results across 15 severity-5 corruptions. Acc., AUC, and
    H-score are percentages. R1 uses KNN feature distance; R2 uses energy
    scoring with percentile calibration; R3 ACI experiments are running; R4
    corrects replay-loss scaling.}
    \label{tab:main-results}
    \centering
    \small
    \setlength{\tabcolsep}{3pt}
    \begin{tabular}{llrrr}
        \hline
        Dataset & Method & Acc. & AUC & H-score \\
        \hline
        CIFAR-10-C & Source & 57.34 & 70.39 & 62.32 \\
        + Noise & Norm-test & 72.95 & 68.51 & \textbf{70.58} \\
        & KNN--STAMP (R1) & \textbf{78.64} & 59.10 & 63.61 \\
        & Energy--STAMP (R2) & \textbf{78.99} & 54.33 & 63.80 \\
        & STAMP+BSN (R4) & 78.95 & \textbf{83.96} & \textbf{81.19} \\
        & Entropy+ACI (R3) & \multicolumn{3}{c}{\emph{running}} \\
        & Energy+ACI (R3) & \multicolumn{3}{c}{\emph{running}} \\
        \hline
        CIFAR-100-C & Source & 35.83 & 43.13 & 38.01 \\
        + Noise & Norm-test & \textbf{45.90} & \textbf{80.91} & \textbf{58.49} \\
        & KNN--STAMP (R1) & 45.73 & 69.93 & 54.15 \\
        & Energy--STAMP (R2) & \textbf{57.14} & \textbf{95.12} & \textbf{71.33} \\
        & STAMP+BSN (R4) & \textbf{58.02} & \textbf{98.29} & \textbf{72.89} \\
        & Entropy+ACI (R3) & \multicolumn{3}{c}{\emph{running}} \\
        & Energy+ACI (R3) & \multicolumn{3}{c}{\emph{running}} \\
        \hline
    \end{tabular}
\end{table}

\subsection{Round-2 Results}
Energy--STAMP produces sharply different outcomes across the two benchmarks.
On CIFAR-10-C, it increases accuracy only slightly relative to KNN--STAMP
(78.99 versus 78.64) while decreasing AUC from 59.10 to 54.33. The H-score
therefore remains low (63.80), 6.78 points below the normalization baseline.
The second intervention consequently does not resolve the CIFAR-10-C
recognition--rejection trade-off.

On CIFAR-100-C, in contrast, energy--STAMP improves on KNN--STAMP by 11.41
accuracy points, 25.19 AUC points, and 17.18 H-score points. It also exceeds
test-time normalization by 11.24 points in accuracy, 14.21 points in AUC, and
12.84 points in H-score. This is the strongest result in Table~\ref{tab:main-results}
and supports the idea that logit energy can provide a more useful memory-admission
signal than a fixed clean feature bank in the 100-class setting. The round-4
correction below further improves these completed results.

The contrasting results prevent a universal claim for energy scoring. The next
analysis should report ID rejection and OOD admission rates for entropy, KNN,
and energy, with per-corruption results, to determine whether the low
CIFAR-10-C AUC arises from score overlap, threshold calibration, or changes in
the samples admitted to memory.

\subsection{Round-4 Results: Replay-Loss Scaling Correction}
STAMP+BSN corrects the replay-loss normalization during memory warm-up. On
CIFAR-10-C, it preserves the energy variant's accuracy (78.95 versus 78.99)
while increasing AUC from 54.33 to 83.96 and H-score from 63.80 to 81.19. It
therefore exceeds test-time normalization by 6.00 accuracy points, 15.45 AUC
points, and 10.61 H-score points. On CIFAR-100-C, STAMP+BSN reaches 58.02
accuracy, 98.29 AUC, and 72.89 H-score, improving on energy--STAMP by 0.88,
3.17, and 1.56 points, respectively. It also exceeds normalization by 12.12
accuracy points, 17.38 AUC points, and 14.40 H-score points.

The large CIFAR-10-C AUC change is consistent with the scale analysis: early,
partially filled replay buffers previously received overly large descent
updates, which could destabilize the model before the memory reached capacity.
However, the aggregate results alone do not establish this mechanism. Per-step
memory occupancy, loss magnitude, and admission rates should be logged in a
follow-up ablation to confirm that the improvement occurs during warm-up rather
than through an uncontrolled experimental difference.

\subsection{Round-3 Experiments in Progress}
The pending round-3 rows evaluate adaptive conformal inference (ACI) as a
threshold controller. Entropy--STAMP+ACI holds STAMP's original score fixed and
tests the threshold intervention alone. Energy--STAMP+ACI tests whether the
strong CIFAR-100-C energy result can be retained while improving the weak
CIFAR-10-C energy AUC. Both variants use the same datasets, corruption order,
batch size, reset policy, and metrics as rounds 1 and 2. When available, their
results will be accompanied by the ACI quantile trajectory and per-corruption
ID-rejection and OOD-admission rates; these diagnostics are necessary because
ACI controls admission calibration rather than directly optimizing AUC.
